%!TEX root main.tex
\section{Infiltration-finite series and infiltration automata}
\label{sec:infiltration}

We have so far considered the Hadamard (\cref{sec:Hadamard}) and the shuffle products (\cref{sec:shuffle}),
motivated by their applications to polyrec sequences~(\cref{sec:polyrec}),
resp., \CDA~power series~(\cref{sec:CDA}).
%
Besides Hadamard and shuffle, in the literature one finds another product,
called \emph{infiltration} (\cref{sec:infiltration:preliminaries}).
%
We observe that the infiltration product falls under the scope of our techniques.
%
In a manner analogous to the Hadamard and shuffle products,
we can define \emph{infiltration-finite series} (\cref{sec:infiltration-finite series})
and corresponding \emph{infiltration automata} (\cref{sec:infiltration automata}),
which are novel models of computation not previously investigated.
%
The main result of this section (\cref{sec:infiltration:equality}) is that
infiltration-finite series constitute an effective prevariety
and thus the commutativity problem is decidable for this class
(\cref{thm:infiltration finite - effective prevariety}).

\subsection{Preliminaries}
\label{sec:infiltration:preliminaries}
\subsubsection{Infiltration algebras}

An \emph{infiltration function} of an algebra $\tuple{A; 0, 1, {+}, {*}}$ (in short, \emph{infiltration})
is a linear function $\delta : A \to A$ satisfying the following product rule
%
\begin{align}
    \label{eq:infiltration rule}
    \tag{infiltration rule}
    \delta(\alpha * \beta) = \delta \alpha * \beta + \alpha * \delta \beta + \delta \alpha * \delta \beta,
    \quad\text{for all } \alpha, \beta \in A.
\end{align}
%
An \emph{infiltration algebra}
$\tuple{A; 0, 1, {+}, {*}, (\delta_j)_{1\leq j\leq d}}$ is an algebra
together with finitely many, not necessarily commuting infiltration functions $\delta_1, \dots, \delta_d : A \to A$.
%
The product rule \cref{eq:infiltration rule} bears resemblance to that for
endomorphisms~\cref{eq:endomorphism} and for derivations~\cref{eq:Leibniz rule}.
%
In the next section we clarify this connection,
which will be exploited to simplify later proofs.

\subsubsection{$\sigma$-differential algebras}

Let $\sigma : A \to A$ be an endomorphism.
A \emph{$\sigma$-derivation} is a linear function $\delta : A \to A$ satisfying 
%
\begin{align}
    \label{eq:sigma-derivation rule}
    \tag{$\sigma$-derivation rule}
    \delta(\alpha * \beta) = \delta \alpha * \beta + \sigma \alpha * \delta \beta,
        \quad\text{for all } \alpha, \beta \in A.
\end{align}
%
If we take the identity endomorphism $\sigma = \id$ then we recover the notion of derivation~\cref{eq:Leibniz rule}.
%
We observe that infiltrations arise as $\sigma$-derivations of a special kind.
% which will shorten some proofs later.
%
\begin{lemma}
    \label{lem:fundamental relationship}
    Let $\sigma, \delta : A \to A$ be two functions satisfying the equation
    %
    \begin{align}
        \label{eq:fundamental relationship}
        \tag{$\sigma$-$\delta$}
        \sigma \alpha = \alpha + \delta \alpha,
            \text{ for every } \alpha \in A.
    \end{align}
    %
    \begin{enumerate}

        \item If $\sigma$ is an endomorphism,
        then the unique $\delta$ satisfying~\cref{eq:fundamental relationship} is an infiltration.
        \\
        We call $\delta$ the infiltration \emph{induced} by $\sigma$.

        \item If $\delta$ is an infiltration,
        then the unique $\sigma$ satisfying~\cref{eq:fundamental relationship} is an endomorphism.
        \\
        We call $\sigma$ the endomorphism \emph{induced} by $\delta$.

    \end{enumerate}
    %
    In both cases, $\delta$ is a $\sigma$-derivation.
\end{lemma}
%
\noindent
For instance, the identity endomorphism $\sigma = \id$
induces the trivial infiltration $\delta \alpha = 0$ (and vice versa).
%
We will use the lemma in the next section to extract an infiltration from an endomorphism,
and conversely to extract an endomorphism from an infiltration
(\cf~the proof of~\cref{lem:infiltration - unique extension}).
%
\begin{proof}
    Regarding point (1), let $\sigma$ be an endomorphism
    and consider the unique map $\delta$ satisfying~\cref{eq:fundamental relationship}.
    Linearity is clear.
    Regarding products, we have
    %
    %\begin{align*}
        $\delta (\alpha \cdot \beta)
        = \sigma (\alpha \cdot \beta) - \alpha \cdot \beta
        = \sigma \alpha \cdot \sigma \beta - \alpha \cdot \beta
        = (\alpha + \delta \alpha) \cdot (\beta + \delta \beta) - \alpha \cdot \beta
        = \alpha \cdot \delta \beta + \delta \alpha \cdot \beta + \delta \alpha \cdot \delta \beta$.
        % = \delta \alpha \cdot \beta + \sigma \alpha \cdot \delta \beta.
    % \end{align*}
    %
    Regarding point (2), let $\delta$ be an infiltration
    and let $\sigma$ be the unique map satisfying~\cref{eq:fundamental relationship}.
    %
    % Clearly $\sigma$ has been chosen so that $\delta$ is a $\sigma$-derivation.
    % It remains to check that $\sigma$ is an infiltration algebra endomorphism.
    Linearity is clear.
    Regarding products we have
    %
    % \begin{align*}
        $\sigma (\alpha \cdot \beta)
            = \alpha \cdot \beta + \delta (\alpha \cdot \beta)
            = \alpha \cdot \beta + \delta \alpha \cdot \beta + \alpha \cdot \delta \beta + \delta \alpha \cdot \delta \beta
            % = \alpha \cdot \beta + \delta \alpha \cdot \beta + \sigma \alpha \cdot \delta \beta
            = \sigma \alpha \cdot \beta + \sigma \alpha \cdot \delta \beta
            % = \sigma \alpha \cdot (\beta + \delta \beta)
            = \sigma \alpha \cdot \sigma \beta$.
    % \end{align*}
\end{proof}

\subsubsection{Infiltration algebras of polynomials}
%
As an example of an infiltration algebra,
consider the algebra of polynomials $\tuple{\poly \Q X; 0, 1, {+}, {\cdot}}$.
%
An infiltration of this algebra is thus a linear function $\Delta : \poly \Q X \to \poly \Q X$
satisfying the product rule
%
\begin{align*}
    \Delta (\alpha \cdot \beta)
        = \Delta \alpha \cdot \beta + \alpha \cdot \Delta \beta + \Delta \alpha \cdot \Delta \beta
        = \Delta \alpha \cdot \beta + S \alpha \cdot \Delta \beta,
\end{align*}
%
where we have defined $S \alpha := \alpha + \Delta \alpha$.
%
By~\cref{lem:fundamental relationship}, the function $\Delta$ is an infiltration iff $S$ is an endomorphism.
Thus the polynomial algebra with an infiltration $\Delta$ carries the structure of an $S$-differential algebra.
%
%For instance, $\Delta (X_1 \cdot X_2) = \Delta X_1 \cdot X_2 + X_1 \cdot \Delta X_2 + \Delta X_1 \cdot \Delta X_2$.
%
We have seen that any function $\Delta : X \to \poly \Q X$
extends uniquely to an endomorphism, resp., a derivation $\Delta : \poly \Q X \to \poly \Q X$.
%
We now state and prove the analogous property for infiltrations,
which will be used later in the definition of infiltration automata.
%
\begin{lemma}
    \label{lem:infiltration - unique extension}
    Any function $\Delta : X \to \poly \Q X$
    extends to a unique infiltration function $\Delta : \poly \Q X \to \poly \Q X$.
\end{lemma}
%
\begin{proof}
    Consider the map $S : X \to \poly \Q X$
    \st~$S X_i := X_i + \Delta X_i$ for all $1 \leq i \leq k$
    and extend it to an endomorphism $S : \poly \Q X \to \poly \Q X$. % as above (by evaluation of polynomials).
    %
    Now extend $\Delta$ to the whole polynomial algebra % $\Delta : \poly \Q X \to \poly \Q X$
    by letting $\Delta \alpha := S \alpha - \alpha$ for all $\alpha \in \poly \Q X$.
    %
    Since $S$ is an endomorphism
    and the pair $S, \Delta$ by definition satisfies~\cref{eq:fundamental relationship},
    by~\cref{lem:fundamental relationship} we have that $\Delta$ is an infiltration function.
    %
    Uniqueness can be shown by induction on the total degree of polynomials.
    %
    We given another argument. %invert all steps in the definition of $\Delta$ above.
    %
    Let $\Delta_1, \Delta_2 : \poly \Q X \to \poly \Q X$ be two infiltration functions
    extending $\Delta : X \to \poly \Q X$.
    %
    Then $S_1, S_2 : \poly \Q X \to \poly \Q X$ defined by
    $S_i \alpha := \alpha + \Delta_i \alpha$ ($\alpha \in \poly \Q X$, $i \in \set{1, 2}$)
    are endomorphisms by~\cref{lem:fundamental relationship}.
    %
    But $S_1, S_2$ agree on $X$, and thus $S_1 = S_2$,
    implying $\alpha + \Delta_1 \alpha = \alpha + \Delta_2 \alpha$ ($\forall \alpha \in \poly \Q X$) and thus $\Delta_1 = \Delta_2$.
\end{proof}
%
% If $\Delta_{a \in \Sigma}$ is a family of infiltrations
% %
% \begin
% \begin{remark}[Differential algebra of polynomials]
%     For every $a \in \Sigma$,
%     consider the mapping $S_a : \poly \Q X \to \poly \Q X$ \st
%     %
%     \begin{align}
        
%         \text{ for all } \alpha \in \poly \Q X.
%     \end{align}
%     %
%     This allows us to write the infiltration product rule as:
%     %

%     %
%     % In short, $S_a := 1 + \Delta_a$.
%     Since $\Delta_a$ is an infiltration function,
%     by~\cref{lem:fundamental relationship} we have that $S_a$ is an endomorphism of the polynomial ring
%     % and thus $\Delta_a$ is an $S_a$-derivation of the polynomial ring for every $a \in \Sigma$,
%     and thus $\tuple{\poly \Q X; {+}, {\cdot}, \tuple{S_a}_{a \in \Sigma}, \tuple{\Delta_a}_{a \in \Sigma}}$ is an $S$-differential algebra.
% \end{remark}

\subsubsection{Infiltration algebras of series}

We show that series can be endowed with the structure of an infiltration algebra.
While for polynomials the notion of product is fixed
and we have used it to define infiltration functions,
for series we proceed the other way around:
The transition structure given by left derivatives is fixed
and we define the infiltration product
by demanding that left derivatives be infiltration functions \wrt~this product.
%
Formally, the \emph{infiltration product} of two series $f, g \in \series \Q \Sigma$,
denoted by $f \infiltration g$,
is defined coinductively as follows:
%
\begin{align}
    \tag{${\infiltration}$-$\e$}
    \label{eq:infiltration:base}
    \coefficient \e {(f \infiltration g)}
        &= f_\e \cdot g_\e, \\
    \tag{${\infiltration}$-$\deriveleft a$}
    \label{eq:infiltration:step}
        \deriveleft a (f \infiltration g)
        &= \deriveleft a f \infiltration g
            + f \infiltration \deriveleft a g
            + \deriveleft a f \infiltration \deriveleft a g,
        \ \forall a \in \Sigma.
\end{align}
%
To see why this uniquely defines a series $f \infiltration g$,
we can reason by induction on the length of words.
%
The first rule gives us the value for the constant term. % $\coefficient \e (f \infiltration g)$.
For nonempty words, we have
%
\begin{align*}
    &\coefficient {a \cdot w} (f \infiltration g)
    = \coefficient w (\deriveleft a (f \infiltration g))
%    = \coefficient w (\deriveleft a f \infiltration g
%        + f \infiltration \deriveleft a g
%        + \deriveleft a f \infiltration \deriveleft a g)
    = \coefficient w (\deriveleft a f \infiltration g)
    + \coefficient w (f \infiltration \deriveleft a g)
    + \coefficient w (\deriveleft a f \infiltration \deriveleft a g),
\end{align*}
%
where the three coefficients for word $w \in \Sigma^*$ are obtained by the inductive assumption.
%
% For instance, a short calculation shows that $ab \infiltration a = 2aab + aba + ab$.
%
Infiltration product is an associative, commutative, and bilinear operation on series,
with identity the series $1 \cdot \e$ (the same identity as for the shuffle product).
%
We call the resulting structure $\series \Q \Sigma_{\infiltration} := \tuple{\series \Q \Sigma; \zero, 1 \cdot \e, {+}, {\infiltration}}$
the \emph{infiltration algebra} of series.
%
%The rule~\cref{eq:infiltration:step} makes $\deriveleft a$ an infiltration of this algebra.
%
Infiltration has been introduced by Chen, Fox, and Lyndon~\cite{ChenFoxLyndon:AM:1958}
in the context of the \emph{free differential calculus} of Fox~\cite{Fox:AM:1953}.
%
More recently, it has been considered in a coalgebraic setting~\cite[Sec.~8]{BasoldHansenPinRutten:MSCS:2017}.
%
From a concurrency perspective, it models the \emph{synchronising interleaving}
where processes are allowed (but not forced) to jointly perform the input actions
(thanks to the term $\deriveleft a f \infiltration \deriveleft a g$ in~\cref{eq:infiltration:step}).
%
For instance, in the interleaving semantics we have $ab \shuffle a = 2aab + aba$,
but in the synchronising interleaving we have $ab \infiltration a = 2aab + aba + ab$.

\begin{remark}
    We have chosen a coinductive approach to the definition of infiltration product
    since it brings to the foreground the similarity with Hadamard and shuffle products.
    %
    Alternatively, one can define the infiltration product on finite words by induction on their length,
    and then lift it to series by linearity and continuity~\cite[page 126]{Lothaire:CUP:1997}.
\end{remark}

\begin{remark}[Differential algebra of series]
    %
    For every input symbol $a \in \Sigma$,
    consider the linear map $\sigma_a : \series \Q \Sigma \to \series \Q \Sigma$ defined by
    %
    \begin{align}
        \tag{$\sigma_a$}
        \sigma_a f := f + \deriveleft a f,
        \text{ for every } f \in \series \Q \Sigma. %, a \in \Sigma.
    \end{align}
    %
    Sometimes we write $\sigma^L_a$ to stress that the definition is \wrt~the left derivative.
    %
    Since $\deriveleft a$ is an infiltration function,
    thanks to~\cref{lem:fundamental relationship}
    we have that $\sigma_a$ is an infiltration algebra endomorphism,
    and thus $\deriveleft a$ is a $\sigma_a$-derivation.
    %
    Using the definition of $\sigma_a$,
    the rule~\cref{eq:infiltration:step} can be rewritten more succinctly as follows:
    %
    \begin{align}
        \tag{${\infiltration}$-$\deriveleft a$'}
        \label{eq:infiltration:step'}
            \deriveleft a (f \infiltration g)
            &= \deriveleft a f \infiltration g
                + \sigma_a f \infiltration \deriveleft a g
            \ \forall a \in \Sigma.
    \end{align}
    %
    This should be compared to the fact that $\deriveleft a$ is a derivation
    (\ie, for the identity endomorphism) of the shuffle algebra.
    %
    % This helps explaining why the infiltration product behaves more similarly to
    % the shuffle product than to the Hadamard one.
\end{remark}

In the next lemma we observe that also right derivatives are infiltration functions.
This shows that infiltration product can be defined equivalently from the left and from the right.
%
This is a symmetry shared by the Hadamard and shuffle products,
however there is no reason why it should hold for an arbitrary product.
% (for instance, it fails for the concatenation product).
%
We will be using this property in the proof of closure under right derivatives for infiltration-finite series
(\cref{lem:infiltration closure}).
%
In the following, define $\sigma^R_a f := f + \deriveright a f$, for all $f \in \series \Q \Sigma$.
%
\begin{lemma}
    \label{lem:derive right infiltration}
    The right derivative operator $\deriveright a$ (with $a \in \Sigma$)
    is an infiltration function of the infiltration algebra:
    it is linear and it satisfies
    %
    \begin{align}
        \label{eq:derive right infiltration}
        \tag{${\infiltration}$-$\deriveright a$}
        \deriveright a (f \infiltration g)
        &= \deriveright a f \infiltration g
            + f \infiltration \deriveright a g
            + \deriveright a f \infiltration \deriveright a g
        = \deriveright a f \infiltration g + \sigma^R_a \infiltration \deriveright a g,
        \quad \forall f, g \in \series \Q \Sigma. % a \in \Sigma.
    \end{align}
\end{lemma}
%
\begin{proof}
    We proceed by coinduction.
    The constant terms of the two sides coincide:
    %
    \begin{align*}
        \coefficient \e (\deriveright a (f \infiltration g))
            % &= \coefficient a (f \infiltration g)
            &= \coefficient \e (\deriveleft a (f \infiltration g))
            = \coefficient \e (\deriveleft a f \infiltration g + f \infiltration \deriveleft a g + \deriveleft a f \infiltration \deriveleft a g)
            = f_a \cdot g_\e + f_\e \cdot g_a + f_a \cdot g_a = \\
            &= \coefficient \e (\deriveright a f \infiltration g + f \infiltration \deriveright a g + \deriveright a f \infiltration \deriveright a g).
    \end{align*}
    %
    Now consider an input symbol $b \in \Sigma$.
    %
    Thanks to~\cref{lem:derive left right commutativity},
    we obtain the following commutation rules:
    %
    \begin{align}
        \label{eq:commutation rules}
        \deriveright a  \deriveleft b f = \deriveleft b \deriveright a f,
            \qquad
        \sigma^R_a \delta^L_b f = \delta^L_b \sigma^R_a f,
            \qquad
        \delta^R_a \sigma^L_b f = \sigma^L_b \delta^R_a f,
            \quad\text{and}\quad
        \sigma^R_a \sigma^L_b f = \sigma^L_b \sigma^R_a f,
        \quad \forall f \in \series \Q \Sigma.
    \end{align}
    %
    % We will write $f^b$ instead of $\deriveleft b f$ for short.
    Thanks to the rules above, we have
    %
    \begin{align*}
        \deriveleft b \deriveright a (f \infiltration g)
            &= \deriveright a \deriveleft b (f \infiltration g) =
                &&\text{by~\cref{eq:commutation rules}} \\
            &= \deriveright a (\deriveleft b f \infiltration g + \sigma^L_b f \infiltration \deriveleft b g) =
                &&\text{(def.~$\infiltration$)} \\
            &= \deriveright a (\deriveleft b f \infiltration g) + \deriveright a (\sigma^L_b f \infiltration \deriveleft b g) =
                &&\text{(linearity of $\deriveright a$)} \\
            &=  \deriveright a \deriveleft b f \infiltration g + \sigma^R_a \deriveleft b f \infiltration \deriveright a g
                + \deriveright a \sigma^L_b f \infiltration \deriveleft b g + \sigma^R_a \sigma^L_b f \infiltration \deriveright a \deriveleft b g =
                &&\text{(coinduction)} \\
            &=  \deriveleft b \deriveright a f \infiltration g + \deriveleft b \sigma^R_a f \infiltration \deriveright a g
                + \sigma^L_b \deriveright a f \infiltration \deriveleft b g + \sigma^L_b \sigma^R_a  f \infiltration \deriveleft b \deriveright a g =
                &&\text{by~\cref{eq:commutation rules}} \\
            &= \deriveleft b (\deriveright a f \infiltration g)
                + \deriveleft b (\sigma^R_a f \infiltration \deriveright a g)
                &&\text{(def.~$\infiltration$)} \\
            &= \deriveleft b (\deriveright a f \infiltration g + \sigma^R_a f \infiltration \deriveright a g).
                &&\text{(linearity of $\deriveleft b$)} \qedhere
    \end{align*}
    %
    % The second equation follows from \cref{lem:fundamental relationship},
    % since we have shown that $\deriveright a$ is an infiltration function
    % and thus $\sigma^R_a$ is its induced endomorphism.
\end{proof}

% lax begin Lax619925.Infiltration
\subsection{Infiltration-finite series}
\label{sec:infiltration-finite series}

We have seen that the Hadamard and shuffle products induce corresponding classes of
Hadamard-finite~(\cref{sec:Hadamard-finite series}), resp.
shuffle-finite series~(\cref{sec:shuffle-finite series}),
together with natural closure properties.
%
The infiltration product admits an analogous development.
%
For series $f_1, \dots, f_k \in \series \Q \Sigma$,
we denote by $\poly \Q {f_1, \dots, f_k}_{\infiltration}$
the infiltration subalgebra finitely generated by the $f_i$'s.
%
% This algebra is usually not free,
% since there may be nontrivial relations between the generators.
%
An infiltration algebra is \emph{differential}
if it is closed under $\deriveleft a$, for all $a \in \Sigma$.

\begin{definition}
    A series is \emph{infiltration finite} if it belongs to a
    finitely generated differential infiltration algebra of series.
\end{definition}
%
The following lemma unpacks the definition above
and provides our working definition for infiltration finite series.
%
Its statement and proof are similar to~\cref{lem:Hadamard finite - working definition}:
We present the former but we omit the latter.
%
\begin{lemma}
    \label{lem:infiltration finite - working definition}
    A series $f \in \series \Q \Sigma$ is infiltration finite
    iff there are generators $f_1, \dots, f_k \in \series \Q \Sigma$ \st:
    \begin{enumerate}
        \item $f \in \poly \Q {f_1, \dots, f_k}_{\infiltration}$ and
        \item $\deriveleft a f_i \in \poly \Q {f_1, \dots, f_k}_{\infiltration}$
        for every $a \in \Sigma$ and $1 \leq i \leq k$.
    \end{enumerate}
    %
    Moreover, we can assume \wlg~that $f = f_1$.
\end{lemma}

Infiltration-finite series have several pleasant closure properties.

\begin{restatable}{lemma}{lemInfiltrationClosure}
    \label{lem:infiltration closure}
    The class of infiltration-finite series
    is an effective differential infiltration algebra:
    It contains $\zero$, $1 \cdot \e$, and it is effectively closed under
    scalar product $c \cdot f$,
    sum $f + g$,
    infiltration product $f \infiltration g$,
    and left derivative $\deriveleft a f$.
    %
    Moreover,
    %
    \begin{itemize}
        \item if $f$ is infiltration finite,
        then $\deriveright a f$ is effectively infiltration finite, and
        %
        \item If $g \in \series \Q \Sigma$ is an anti-derivative
        of a tuple of infiltration-finite series $f = \tuple{f_a}_{a \in \Sigma} \in \series \Q \Sigma^\Sigma$,
        then $g$ is effectively infiltration finite.
    \end{itemize}
\end{restatable}
%
\begin{proof}
    All closure properties are based on simple manipulations of the generators.
    They are based on~\cref{lem:infiltration finite - working definition}
    and follow the same pattern as for Hadamard-finite and shuffle-finite series
    (\cf~\cref{lem:Hadamard closure properties}).
    %
    As an illustration, we show closure under right derivatives.
    %
    Let $f$ be an infiltration-finite series and fix an input symbol $a \in \Sigma$.
    %
    By~\cref{lem:infiltration finite - working definition},
    there are generators $f_1, \dots, f_k$ with $f = f_1$ generating the $\deriveleft {}$-closed infiltration algebra
    %
    % \begin{align*}
        $A := \poly \Q {f_1, \dots, f_k}_{\infiltration}$.
    % \end{align*}
    %
    We adjoin new generators $\deriveright a f_1, \dots, \deriveright a f_k$,
    obtaining the larger infiltration algebra
    %
    % \begin{align*}
        $B := \poly \Q {f_1, \dots, f_k, \deriveright a f_1, \dots, \deriveright a f_k}_{\infiltration}$.
    % \end{align*}
    %
    By construction, $\deriveright a f = \deriveright a f_1$ is in $B$.
    It remains to show that $B$ is $\deriveleft {}$-closed.
    %
    Since $A$ is $\deriveleft {}$-closed, it suffices to show this for the new generators.
    %
    So consider $\deriveleft b \deriveright a f_i$
    for arbitrary $a, b \in \Sigma$ and $1 \leq i \leq k$.
    %
    By~\cref{lem:derive left right commutativity}, left and right derivatives commute,
    so we have $\deriveleft b \deriveright a f_i = \deriveright a \deriveleft b f_i$.
    %
    Since $A$ is $\deriveleft {}$-closed, $\deriveleft b f_i$ is in $A$,
    and thus we can write $\deriveleft b f_i = p_i^b(f_1, \dots, f_k)$
    for some polynomial $p_i^a$ (where product is interpreted as infiltration product).
    %
    Thanks to linearity of right derivative and~\cref{lem:derive right infiltration},
    a proof by structural induction on polynomials
    shows that $\deriveright a (p_i^a(f_1, \dots, f_k))$ can be written as a polynomial expression
    $q_i^{a, b}(f_1, \dots, f_k, \deriveright a f_1, \dots, \deriveright a f_k)$.
    %
    We thus have
    %
    $\deriveleft b \deriveright a f_i
            = \deriveright a \deriveleft b f_i
            = q_i^{a, b}(f_1, \dots, f_k, \deriveright a f_1, \dots, \deriveright a f_k)
            \in B$,
    as required.
\end{proof}

In this section we have defined the class of infiltration-finite series in a semantic fashion,
based on which we have shown effective closure properties~(\cref{lem:infiltration closure}).
%
In order to show that infiltration-finite series constitute an effective prevariety,
it remains to show that the zeroness problem is decidable for them.
%
To this end, we introduce in the next section a suitable automaton model capturing the same class.

\subsection{Infiltration automata}
\label{sec:infiltration automata}

The syntax and semantics of infiltration automata is largely parallel to that of Hadamard and shuffle automata.
An \emph{infiltration automaton} is a tuple
$\AA = \tuple{\Sigma, X, F, \Delta}$
where $\Sigma$ is a finite \emph{input alphabet},
$X = \set{X_1,  \dots X_k}$ is a finite set of \emph{nonterminal symbols},
$F : X \to \Q$ is the \emph{output function},
and $\Delta : \Sigma \to X \to \poly \Q X$ is the \emph{transition function}.
%
The only difference is how to extend the transition function $\Delta_a : X \to \poly \Q X$
from variables $X$ to configurations $\poly \Q X$.
%
Here, we extend $\Delta_a$ to a unique infiltration $\Delta_a : \poly \Q X \to \poly \Q X$,
thanks to~\cref{lem:infiltration - unique extension}.
%
In a way entirely parallel to Hadamard and shuffle automata,
this defines a semantic function $\Asem \AA \_ : \poly \Q X \to \series \Q \Sigma$.
%
When the automaton $\AA$ is fixed by the context, we also just write $\sem \_$.
%
The main property of the semantics is that it is a homomorphism,
in complete analogy for the corresponding properties for Hadamard and shuffle automata
(\cref{lem:Hadamard automata - properties of the semantics},
resp., \cref{lem:shuffle automata - properties of the semantics}).
%
\begin{lemma}[Properties of the semantics]
    \label{lem:infiltration automata - properties of the semantics}
    The semantic function $\Asem \AA \_$ is a homomorphism
    from the $S$-differential algebra of polynomials
    to the $\sigma$-differential infiltration algebra of series.
    %
    In other words, for every configurations $\alpha, \beta \in \poly \Q X$,
    %
    \begin{align}
        \label{eq:infiltration:sem:scalar-product}
        \Asem \AA {c \cdot \alpha}
            &= c \cdot \Asem \AA \alpha,
            &&\forall c \in \Q, \\
        \label{eq:infiltration:sem:sum}
        \Asem \AA {\alpha + \beta}
            &= \Asem \AA \alpha + \Asem \AA \beta, \\
        \label{eq:infiltration:sem:product}
        \Asem \AA {\alpha \cdot \beta}
            &= \Asem \AA \alpha \infiltration \Asem \AA \beta, \\
        \label{eq:infiltration:sem:endomorphism}
        \Asem \AA {S_a \alpha}
            &= \sigma_a (\Asem \AA \alpha),
            &&\forall a \in \Sigma \\
        \label{eq:infiltration:sem:derivation}
        \Asem \AA {\Delta_a \alpha}
            &= \deriveleft a (\Asem \AA \alpha),
            &&\forall a \in \Sigma.
        \end{align}
\end{lemma}
%
\begin{proof}
    Property~\cref{eq:infiltration:sem:derivation}
    holds by definition of the semantics.
    %
    Properties~\cref{eq:infiltration:sem:scalar-product} and~\cref{eq:infiltration:sem:sum}
    follow from linearity of $\deriveleft a$ and $\Delta_a$.
    %
    From these, \cref{eq:infiltration:sem:endomorphism} follows easily:
    %
    $
        \sem {S_a \alpha}
        = \sem {\alpha + \Delta_a \alpha}
        = \sem \alpha + \sem {\Delta_a \alpha}
        = \sem \alpha + \deriveleft a \sem \alpha
        = \sigma_a \sem \alpha$.
    %
    For the last property~\cref{eq:infiltration:sem:product}, we proceed coinductively.
    %
    Clearly, the constant term of $\sem {\alpha \cdot \beta}$ and $\sem \alpha \infiltration \sem \beta$
    is the same, namely $F \alpha \cdot F \beta$.
    %
    For the coinductive step, we show that the left derivatives
    of $\sem {\alpha \cdot \beta}$ and $\sem \alpha \infiltration \sem \beta$ coincide:
    %
    % For readability, we write $\alpha^a$ for $\Delta_a \alpha$ when $\alpha \in \poly \Q X$
    % and $f^a$ for $\deriveleft a f$ when $f \in \series \Q \Sigma$.
    %
    \begin{align*}
        \deriveleft a \sem {\alpha \cdot \beta}
            &= \sem {\Delta_a (\alpha \cdot \beta)}
            = \sem {\Delta_a \alpha \cdot \beta + S_a \alpha \cdot \Delta_a \beta}
            = \sem {\Delta_a \alpha \cdot \beta} + \sem {S_a \alpha \cdot \Delta_a \beta}
            \stackrel ! = \sem {\Delta_a \alpha} \infiltration \sem \beta
                + \sem {S_a \alpha} \infiltration \sem {\Delta_a \beta} = \\
            &= (\deriveleft a {\sem \alpha}) \infiltration \sem \beta
                + (\sigma_a \sem \alpha )\infiltration \deriveleft a \sem \beta
            = \deriveleft a (\sem \alpha \infiltration \sem \beta). \qedhere
    \end{align*}
\end{proof}

The following characterisation connects infiltration automata and infiltration-finite series.
Thanks to it, we can use algorithmic results on infiltration automata for infiltration-finite series.
%
We omit the proof, which is analogous to that of~\cref{lem:Hadamard automata and series},
relying on~\cref{lem:infiltration automata - properties of the semantics,lem:infiltration finite - working definition}.
%
\begin{restatable}{lemma}{lemInfiltrationAutomataInfiltrationFinite}
    \label{lem:infiltration automata = infiltration finite}
    A series is recognised by an infiltration automaton
    if, and only if, it is infiltration finite.
\end{restatable}

\subsection{Equality problem}
\label{sec:infiltration:equality}

We show that the equality problem for infiltration-finite series is decidable.
%
% When $\Sigma = \set a$ (i.e., for univariate sequences, \aka~streams, $\N \to \Q$)
% decidability follows from \cite[Theorem 4.2]{BorealeCollodiGorla:ACMTCL:2024}.
% %
% We show that this can be generalised to arbitrary finite alphabets $\Sigma$.
%
To this end, we leverage a classic technique based on chains of polynomial ideals and Hilbert's finite basis theorem.
%
This technique has been used by Novikov and Yakovenko to derive decidability results
on continuous time and discrete time dynamical systems~\cite{NovikovYakovenko:1999}.
%
In the terminology of this paper,
their results yield decidability of equality for univariate ($\Sigma = \set a$) Hadamard-finite and shuffle-finite series.
%
It has provided the basis of several subsequent decidability results,
such as in the work of Benedikt \etal~for multivariate ($\card \Sigma \geq 1$) Hadamard-finite series
(\aka~polynomial automata \cite{BenediktDuffSharadWorrell:PolyAut:LICS:2017}),
in our work on multivariate shuffle-finite series~\cite{Clemente:CONCUR:2024},
and in the work of Boreale \etal~for univariate infiltration-finite series~\cite{BorealeGorla:CONCUR:2021,BorealeCollodiGorla:ACMTCL:2024}.

Fix an infiltration automaton $\AA = \tuple{\Sigma, X, F, \Delta}$
recognising an infiltration-finite series $f = \Asem \AA {X_1} \in \series \Q \Sigma$.
%
Since equality reduces to zeroness, we show how to decide $\Asem \AA {X_1} = \zero$.
%
For a set of configurations $P \subseteq \poly \Q X$,
let the \emph{polynomial ideal generated by $P$}, written $\ideal P$,
be the set of all polynomials of the form $\beta_1 \cdot \alpha_1 + \cdots + \beta_n \cdot \alpha_n$,
with $\alpha_1, \dots, \alpha_n \in P$ configurations in $P$
and $\beta_1, \dots, \beta_n \in \poly \Q X$ arbitrary configurations.
%
(We refer to \cite{CoxLittleOShea:Ideals:2015} for more details on algebraic geometry.)
%
We will be looking at the ideal $\ideal C$ generated by all reachable configurations
$C := \setof{\Delta_w X_1 \in \poly \Q X}{w \in \Sigma^*}$.
%
The intuition is that, while the set of reachable configurations $C$ is complicated
(formally, its membership problem is undecidable),
the ideal it generates $\ideal C$ is simple:
It is finitely generated (as an ideal)
and we can compute a finite set of generators for it.
%
In other words, $\ideal C$ is effectively finitely generated.
%
Moreover, the ideal abstraction is sound for the zeroness problem:
when the generators vanish, so do all polynomials in the ideal they generate,
and so no information is lost.

For a length $n \in \N$,
consider the polynomial ideal $I_n$
generated by all configurations reachable from $X_1$
by reading words of length $\leq n$:
%
\begin{align*}
    I_n := \idealof {\Delta_w X_1} {w \in \Sigma^{\leq n}}.
\end{align*}
%
The following lemma describes the relevant properties of these ideals.
%
\begin{lemma}
    \label{lem:infiltration ideals}
    \begin{enumerate}
        \item $I_n \subseteq I_{n+1}$.
        \item $\Delta_a I_n \subseteq I_{n+1}$.
        \item $I_{n+1} = I_n + \idealof {\Delta_a I_n} {a \in \Sigma}$.
        \item $I_n = I_{n+1}$ implies $I_n = I_{n+1} = I_{n+2} = \cdots$.
    \end{enumerate}
\end{lemma}
\begin{proof}
    The first point holds since the operation of constructing the ideal generated by a set
    is monotonic for inclusion of sets.
    %
    To prove the second point we use the fact that $\Delta_a$ is an infiltration
    (this is the only property relying directly on this).
    %
    Let $\alpha \in \Delta_a I_n$. Then $\alpha$ is of the form
    %
    \begin{align*}
        \alpha
        &= \Delta_a \sum_{i = 1}^m \beta_i \cdot \Delta_{w_i} X_1
        = \sum_{i = 1}^m \Delta_a (\beta_i \cdot \Delta_{w_i} X_1)
        = \sum_{i = 1}^m (\Delta_a \beta_i \cdot \underbrace{\Delta_{w_i} X_1}_{\in I_n}
            + S_a \beta_i \cdot \underbrace{\Delta_a \Delta_{w_i} X_1}_{\in I_{n+1}}),
    \end{align*}
    %
    which is clearly in $I_{n+1}$.
    %
    We now consider the third point.
    The ``$\supseteq$'' inclusion holds by the first two points
    and the fact that $I_{n+1}$ is an ideal.
    %
    For the other inclusion, let $\alpha \in I_{n+1}$.
    We can write $\alpha$ separating words of length exactly $n+1$ from the rest, and we have
    %
    \begin{align*}
        \alpha
            &= \underbrace \beta_{\in I_n}
            + \sum_{w \in \Sigma^{n+1}} \beta_w \cdot \Delta_w X_1
            = \beta + \sum_{a \in \Sigma} \sum_{w \in \Sigma^n} \beta_{w \cdot a} \cdot \Delta_{w \cdot a} X_1
            = \beta + \sum_{a \in \Sigma} \sum_{w \in \Sigma^n} \beta_{w \cdot a} \cdot \Delta_a \underbrace{\Delta_w X_1}_{\in I_n},
    \end{align*}
    %
    where the latter quantity is clearly in the ideal generated by $\Delta_a I_n$.
    %
    The last point follows from the previous one
    since we have shown that $I_{n+1}$ is a function of $I_n$.
\end{proof}

Now consider the chain of ideals
%
\begin{align}
    \label{eq:infiltration ideal chain}
    I_0 \subseteq I_1 \subseteq \cdots \subseteq \poly \Q X.
\end{align}
%
By \emph{Hilbert's finite basis theorem}~\cite[Theorem 4, §5, Ch. 2]{CoxLittleOShea:Ideals:2015},
there is $N \in \N$ \st~the chain stabilises $I_N = I_{N+1} = \cdots$.
%
Thanks to the last point of~\cref{lem:infiltration ideals}, such an $N$ is computable:
Indeed, $N$ can be taken to be the smallest $n \in \N$ \st~$I_n = I_{n+1}$,
and ideal equality can be decided by checking that the generators of $I_{n+1}$ are all in $I_n$.
% lax begin Lax619925.IdealMembership
The latter test is an instance of the \emph{ideal membership problem},
which is decidable (in exponential space)~\cite{Mayr:STACS:1989}.
% lax end
%
We call such minimal $N$ the \emph{stabilisation index} of the ideal chain.
%
We remark that it is crucial that such an $N$ not only exists
(as follows already from Hilbert's theorem, whose proof is notoriously non-constructive),
but it can be computed.
%
The following property provides a decidable criterion for zeroness.
%
\begin{lemma}
    \label{lem:infiltration zeroness characterisation}
    Let $N \in \N$ be the stabilisation index
    of the ideal chain~\cref{eq:infiltration ideal chain}.
    %
    We then have
    \begin{align*}
        \Asem \AA {X_1} = \zero
        \qquad \text{if, and only if,} \qquad
        \coefficient w (\Asem \AA {X_1}) = 0
        \text{ for all } w \in \Sigma^{\leq N}.
    \end{align*}
\end{lemma}
\begin{proof}
    The ``only if'' direction is trivial.
    %
    For the other direction, assume that the semantics is zero for words of length $\leq N$.
    %
    Thanks to the properties of the semantics (\cref{lem:infiltration automata - properties of the semantics}),
    for every word $w \in \Sigma^*$ we can write
    %
    \begin{align*}
        \deriveleft w \sem {X_1}
        &= \sem {\Delta_w X_1}
        = \sem {\sum_{u \in \Sigma^{\leq N}} \beta_u \cdot \Delta_u X_1}
        = \sum_{u \in \Sigma^{\leq N}} \sem {\beta_u} \infiltration \sem {\Delta_u X_1},
    \end{align*}
    %
    where we have used the fact that $\Delta_w X_1$ is in $I_N$.
    %
    By taking constant terms on both sides, we have
    \begin{align*}
        \coefficient w \sem {X_1}
        &= \coefficient \e (\deriveleft w \sem {X_1})
        = \sum_{u \in \Sigma^{\leq N}} \coefficient \e \sem {\beta_u} \cdot \coefficient \e \sem {\Delta_u X_1}
        = \sum_{u \in \Sigma^{\leq N}} \coefficient \e \sem {\beta_u} \cdot \coefficient \e (\deriveleft u \sem {X_1})
        = \sum_{u \in \Sigma^{\leq N}} \coefficient \e \sem {\beta_u} \cdot \coefficient u \sem {X_1}.
    \end{align*}
    %
    But $u$ has length $\leq N$, and thus by the assumption $\coefficient u \sem {X_1} = 0$,
    implying $\coefficient w \sem {X_1} = 0$, as required.
\end{proof}
%
\begin{restatable}{theorem}{thmEqInfiltrationDec}
    \label{thm:infiltration finite - decidable equality}
    The equality problem for infiltration-finite series is decidable.
\end{restatable}
\begin{proof}%[Proof of~\cref{thm:infiltration finite - decidable equality}]
    By using the effective closure properties of infiltration-finite series,
    we reduce equality $f = g$ to zeroness $f - g = \zero$.
    %
    So let $f = \Asem \AA {X_1}$ be an infiltration-finite series,
    recognised by an infiltration automaton $\AA$.
    %
    Compute the stabilisation index $N \in \N$ of the ideal chain~\cref{eq:infiltration ideal chain}.
    Thanks to~\cref{lem:infiltration zeroness characterisation},
    we can decide $f = \zero$
    by enumerating all words $w$ of length $\leq N$
    and checking $\coefficient w f = 0$.
    %
    The latter test is performed by applying the semantics of the infiltration automaton.
\end{proof}

\begin{remark}[Effective algebraic variety of initial conditions]
    \label{rem:effective algebraic varieties}
    The zeroness algorithm above solves a more general problem.
    %
    Not only it decides whether a given initial condition gives rise to the zero semantics,
    but it also shows that the set of all such initial conditions
    forms an \emph{effective algebraic variety} (in the sense of algebraic geometry)
    and computes a finite representation for it in the form of a finite set of polynomials, as we now explain.
    %
    In fact, the same remark applies to equivalence and commutativity,
    and to Hadamard and shuffle automata.

    Given any set of polynomials $P \subseteq \poly \Q X$,
    let their \emph{zero set} $V(P) \subseteq \C^k$ be the set of their common zeroes,
    %
    \begin{align*}
        V(P) := \setof{x \in \C^k}{\forall p \in P : p(x) = 0}.
    \end{align*}
    %
    Sets of the form $V(P) \subseteq \C^k$ for some $P \subseteq \poly \Q X$
    are called \emph{algebraic varieties}~\cite[§2, Ch.~1]{CoxLittleOShea:Ideals:2015}
    and play a prominent role in algebraic geometry.
    %
    Fix a Hadamard, shuffle, or infiltration automaton $\AA_F = \tuple{\Sigma, X, F, \Delta}$,
    where we have made explicit the dependency of the automaton on the initial condition $F$.
    %
    In the algebraic framework,
    it is more natural to consider a slightly more general class of initial conditions $F : X \to \C$,
    mapping nonterminals to complex numbers.
    %
    For every initial configuration $\alpha \in \poly \Q X$,
    let $\FF_\alpha$ be the set of initial conditions giving rise to the zero semantics from $\alpha$,
    %
    \begin{align*}
        % \FF := \setof{F : X \to \C}{\Asem {\AA_F} {\alpha} \text{ is commutative}}.
        \FF_\alpha := \setof{F : X \to \C}{\Asem {\AA_F} \alpha = \zero}.
    \end{align*}
    %
    Thus, by definition $\FF_\alpha$ is the set of common zeroes $\FF_\alpha = V(P)$ of the (possibly infinite) set of polynomials
    %
    \begin{align*}
        %\label{eq:commutativity polynomials}
        %P := \setof {\Delta_u \alpha - \Delta_v \alpha \in \poly \Q X}{u, v \in \Sigma^*, u \sim v},
        P := \setof {\Delta_w \alpha \in \poly \Q X}{w \in \Sigma^*},
    \end{align*}
    %
    %where recall that $u \sim v$ means that $u, v$ are permutation equivalent.
    %
    and thus $\FF_\alpha$ is an algebraic variety.
    %
    The zeroness problem amounts to decide whether a given initial condition $F$ is in this algebraic variety,
    which is just an instance of the membership problem.
    %
    In fact, we can replace $P$ by the polynomial ideal it generates,
    and obtain that $\FF_\alpha = V(P) = V(\ideal P)$ is also the set of common zeroes of $\ideal P$.
    %
    But $\ideal P$ is just the limit $\bigcup_{n = 0}^\infty I_n$ of the polynomial ideal chain~\cref{eq:infiltration ideal chain}.
    %
    We know that there is a computable $N \in \N$ \st~$\ideal P = I_N = \ideal {P_N}$,
    and thus $\FF_\alpha$ is the set of common zeroes $\FF_\alpha = V(P_N)$
    of the \emph{finite} and \emph{computable} set of polynomials $P_N \subseteq P$
    obtained by considering words of length $\leq N$,
    %
    \begin{align*}
        P_N := \setof {\Delta_w \alpha \in \poly \Q X}{w \in \Sigma^{\leq N}}.
    \end{align*}
    %
    As a consequence, $\FF_\alpha$ is an \emph{effective} algebraic variety,
    since its membership problem $F \in \FF_\alpha$ is decidable:
    it suffices to check whether all polynomials in $P_N$ vanish on $F$.
    %
    We can use $P_N$ to solve the following two variants of the zeroness problem:
    %
    \begin{enumerate}
        \item %{$\exists$-commutativity}: 
        Decide whether \emph{there exists} an initial condition $F : X \to \C$
        giving rise to the zero semantics.

        Thanks to \emph{Hilbert's Nullstellensatz}~\cite[Theorem 1, §1, Ch.~4]{CoxLittleOShea:Ideals:2015},
        this amounts to check whether the ideal generated by $P_N$ does not contain $1$.
        %
        The latter can be established with standard techniques
        from computational algebraic geometry,
        such as \emph{Gröbner bases}~\cite[Ch.~2]{CoxLittleOShea:Ideals:2015}.

        \item %{$\forall$-commutativity}:
        Decide whether \emph{every} initial condition $F : X \to \C$
        gives rise to the zero semantics.
        %
        This is equivalent to $P_N = \set 0$.
    \end{enumerate}
\end{remark}

The following theorem summarises the results on infiltration-finite series that we have obtained.
%
\begin{theorem}
    \label{thm:infiltration finite - effective prevariety}
    The class of infiltration-finite series is an effective prevariety.
    Consequently, the commutativity problem is decidable.
    Finally, the commutativity problem is at least as hard as the equivalence problem (under polynomial time reductions)
\end{theorem}
%
\begin{proof}
    The effective closure properties~\cref{prevariety:A,prevariety:B} 
    have been established in~\cref{lem:infiltration closure}
    and decidability of equality in~\cref{thm:infiltration finite - decidable equality}.
    %
    Decidability of commutativity follows from \cref{thm:commutativity for effective prevarieties}.
    %
    Hardness of commutativity follows from~\cref{lem:equality reduces to commutativity}
    and the fact that infiltration-finite series are effectively closed under anti-derivatives~(\cref{lem:infiltration closure}).
\end{proof}
% lax end